\section{MLOps}

% NOTA: a la fecha de redacción, el repositorio (develop y main) no contiene código, workflows ni
% capturas de MLOps. Los textos marcados con \pendiente{} deben completarse con la evidencia real.

\subsection{Importancia de MLOps y CRISP-ML(Q)}

Un modelo de machine learning no es un artefacto estático como un binario: su comportamiento depende de los datos con los que se entrenó, y esos datos cambian con el tiempo. MLOps aplica las prácticas de DevOps (versionado, automatización, pruebas, despliegue continuo y monitoreo) a todo el ciclo de vida del modelo: datos, entrenamiento, validación, despliegue y observación en producción. Sin MLOps, los modelos se entrenan a mano en notebooks, no se pueden reproducir, se despliegan sin validar y se degradan sin que nadie lo note (\textit{data drift} o \textit{concept drift}).

CRISP-ML(Q) (\textit{Cross-Industry Standard Process for Machine Learning with Quality assurance}) extiende el clásico CRISP-DM con aseguramiento de calidad en cada fase. Define seis fases iterativas:
\begin{enumerate}
    \item Entendimiento del negocio y de los datos.
    \item Preparación de datos.
    \item Modelado.
    \item Evaluación.
    \item Despliegue.
    \item Monitoreo y mantenimiento.
\end{enumerate}
En cada fase se identifican los riesgos que podrían afectar la calidad del modelo y se definen métodos para mitigarlos. Así cada etapa tiene criterios de salida verificables, lo que encaja con los gates del pipeline ya usados en DevSecOps.

En MisiOps, la constitución del proyecto define como caso de uso de MLOps \textbf{entrenar modelos sobre el historial de transacciones del usuario para predecir su ahorro a fin de mes}. Esa predicción responde directamente al problema de negocio: que el usuario sepa con anticipación si llegará a fin de mes con ahorro o con deuda.

\subsection{Arquitectura de MLOps y componentes}

\pendiente{diagrama de arquitectura de MLOps (\texttt{figures/mlops/arquitectura.png}) y explicación de cada componente: origen de datos, \textit{feature store} o \textit{pipeline} de features, registro de experimentos y modelos, servicio de inferencia y monitoreo.}

\subsection{Pipeline de CI/CD en GitHub Actions}

\pendiente{workflow de MLOps (archivo en \texttt{.github/workflows/}), sus etapas, sus disparadores y una captura de una corrida exitosa.}

\subsection{Implementación}

\subsubsection{Entendimiento del negocio}

El objetivo es estimar, a mitad de mes, el ahorro con el que el usuario cerrará el periodo presupuestario: ingresos menos gastos. Con esa estimación MisiOps podría alertar con anticipación cuando el ritmo de gasto lleve a cerrar el mes en negativo. \pendiente{métrica de éxito del negocio y métrica técnica (por ejemplo, MAE en soles) con su umbral de aceptación.}

\subsubsection{Feature engineering}

\pendiente{origen y volumen de los datos, variables construidas (por ejemplo, gasto acumulado por categoría, ingresos del periodo, días transcurridos, gasto del mismo periodo del mes anterior) y validaciones de calidad de datos.}

\subsubsection{Selección del modelo}

\pendiente{modelos comparados, métricas obtenidas y criterio de selección del modelo final.}

\subsubsection{Smoke test}

\pendiente{prueba automática que valida el modelo empaquetado: carga, predicción sobre un caso conocido y tiempo de respuesta.}

\subsubsection{Despliegue}

\pendiente{forma de servir el modelo (endpoint, contenedor o batch) y relación con la infraestructura de la sección de IaC.}

\subsubsection{Monitoreo}

\pendiente{métricas que se monitorean en producción (deriva de datos, error de predicción, latencia) y criterio para reentrenar.}

\subsection{Evidencias}

\pendiente{capturas de las corridas del pipeline de MLOps, del registro de experimentos o modelos y del modelo desplegado.}
